\expandafter\def\csname HistoricalNoteText550\endcsname{Zuweilen (obwohl äusserst selten) sind uns allerdings Fälle vorgekommen, die eine Ausnahme davon machten, und wo Spuren von Beschleunigung oder Retardation der Änderung in so kurzen Zwischenzeiten sich doch unverkennbar nachweisen liessen. Mit Ausführlichkeit soll dieser Gegenstand in Zukunft abgehandelt werden,}
\SourceLeaf{542}als die Kette offen ist, fängt gleichfalls in dem Augenblick an mitzuschwingen, wo die Kette geschlossen wird, allein diese Schwingungen, von derselben Dauer, wie die natürlichen, nehmen an Grösse beständig zu, bis diese (erst nach sehr langer Zeit) zu einem Maximum kommt, wo der Widerstand der Luft der Vergrösserung durch die Inductionskraft das Gleichgewicht hält. Um diesen merkwürdigen Versuch wirklich anstellen zu können, wurde (da die Aufhängung eines grossen Stabes wegen Mangel eines zweiten dafür passenden Multiplicators jetzt nicht thunlich war) der einpfündige Stab des physicalischen Cabinets durch Verbindung mit einem ähnlichen etwas schwächer magnetisirten auf bekannte Weise astatisch gemacht, oder vielmehr zu einer Doppelnadel, deren natürliche Schwingungsdauer genau auf $42''3$ gebracht wurde. Der Versuch gelang damit auf das vollkommenste. Der in der Sternwarte schwingende Stab theilte dieser Doppelnadel im physicalischen Cabinet, in dem Augenblick wo die Kette geschlossen wurde, wie durch eine wunderbare Sympathie seine Schwingungen mit, und zwar so, dass jede folgende etwa 50 Scalentheile oder einen halben Grad grösser wurde, als die vorhergehende. Bald ging das ganze Scalenbild aus dem Felde, allein fortwährend konnte man an der immer wachsenden Schnelligkeit, mit welcher das Scalenbild durch das Gesichtsfeld ging, die Zunahme des Schwingungsbogens erkennen. Über eine Stunde wurde dies wunderbar sympathetische Spiel beobachtet.

Es braucht kaum bemerkt zu werden, dass auch der vierpfündige Stab im magnetischen Observatorium in die geschlossene Kette einen Strom inducirt, dessen Dasein an der schnellen Abnahme des Schwingungsbogens auf das bestimmteste erkannt wird, und der daher auch auf die beiden andern Stäbe Wirkungen ausüben muss, denen ähnlich, welche der erstere Versuch gezeigt hat; allein die Rechnung ergibt, und die Erfahrung bestätigt, dass diese Wirkungen zu klein ausfallen, um merklich zu sein. Noch weniger könnte also der schwächste Stab unter den dreien merkliche Wirkungen dieser Art erzeugen.

\clearpage
\SourceLeaf{546}\Needspace{31\baselineskip}\vspace{6mm}
\pdfbookmark[1]{Brief von Gauss an Schumacher. 1836}{letter546}
\begin{center}
{\large\itshape Brief von Gauss an Schumacher.\par}
\vspace{5mm}\hrule\vspace{2mm}
{\small \textsc{Schumacher}. Astronomische Nachrichten Nr.\ 310. 1836. Juni 11.\par}
\vspace{2mm}\hrule\vspace{5mm}
\end{center}
\begin{flushright}Göttingen. 1836. April 23.\end{flushright}
Es waren heute Morgen ausserordentliche Bewegungen der Magnetnadel, noch grösser als am 7. Februar 1835. Dies veranlasste mich einige Sets in der Sternwarte zu beobachten, während Dr.\ \textsc{Goldschmidt} im magnetischen Observa\SourceLeaf{547}torium aufzeichnete. Der gleichförmige Gang bestätigte sich hier so schön, dass ich es wagte den gegenseitigen Uhrstand daraus abzuleiten.

\Needspace{9\baselineskip}
\begin{center}\small
\begin{tabular}{@{}p{0.82\linewidth}r@{}}
Es fand sich, aus einem schnellen Aufsteigen Campa vor Shelton & $4'\,41''1$\\
Aus einem wenige Minuten nachher erfolgten Niedersteigen & $4\quad42.4$\\\cline{2-2}
\multicolumn{1}{r}{Mittel} & $4'\,41''7$
\end{tabular}
\end{center}
Eine directe Vergleichung der Uhren gab,
\begin{center}\small
\begin{tabular}{@{}lll@{}}
1) durch ein Zeichen am Fenster & . . . . . . . . . . & $4'\,41''5$\\
2) durch einen Inductionsimpuls & . . . . . . . . . . & $4\quad41.5$
\end{tabular}
\end{center}
Also eine herrliche Bestätigung dessen, was ich Astronomische Nachrichten Nr.\ 276 [S.\ 539 d.\ B.] gesagt habe.

\clearpage
\pdfbookmark[1]{Das in den Beobachtungsterminen anzuwendende Verfahren.}{verfahren547}
\begin{center}
{\large\itshape Das in den Beobachtungsterminen anzuwendende Verfahren.\par}
\vspace{5mm}\hrule\vspace{2mm}
{\small Resultate aus den Beobachtungen des magnetischen Vereins. 1836. II.\par}
\vspace{2mm}\hrule\vspace{5mm}
\end{center}
Die sechs jährlich festgesetzten Termine fallen gegen das Ende der Monate Januar, März, Mai, Julius, September, November; sie fangen an am letzten Sonnabend in jedem dieser Monate, Mittags nach Göttinger mittlerer Zeit, und schliessen am Mittag des folgenden Tages; die bisher jedem Haupttermine hinzugefügten Nebentermine (Abends von 8—10 Uhr am Dinstag und Mittwoch der folgenden Woche) werden künftig wegfallen.

In jedem Termine wird, der Regel nach, der Stand der Magnetnadel von fünf zu fünf Minuten bestimmt, so dass ein Termin 289 Resultate gibt. In Göttingen wird die Uhr vor Anfang jedes Termins genau auf mittlere Zeit gestellt. Da eine nahe Gleichzeitigkeit der einzelnen Bestimmungen an den verschiedenen Beobachtungsorten sehr wünschenswerth ist, so haben die Beobachter an den meisten andern Orten die Gewohnheit, ihre Uhren gleichfalls auf Göttinger mittlere Zeit zu stellen. Wo dies nicht wohl geschehen kann, ist zu empfehlen, dass man zu den Beobachtungsmomenten diejenigen vollen Minuten der Uhr wähle, die den Göttinger Beobachtungszeiten am nächsten kommen. Hätte man z.\ B.\ vor
\SourceLeaf{548}Anfang des Termins ausgemittelt, dass die bei der Beobachtung zu gebrauchende Uhr um $13'\,48''$ vor Göttinger mittlerer Zeit voraus sei, so würden die Bestimmungen des Standes der Nadel für die Uhrzeiten $0^{\mathrm h}\,14'\,.....\,0^{\mathrm h}\,19'\,.....\,0^{\mathrm h}\,24'\,.....\,0^{\mathrm h}\,29'$ u.\ s.\ f.\ zu machen sein. \textit{Volle} Minuten zu wählen, ist aber jedenfalls anzurathen, weil man sich so die einzelnen Operationen leichter mechanisch macht.

Unter dem Stand der Magnetnadel, welcher für die einzelnen Zeitmomente bestimmt werden soll, ist hier nicht diejenige Stellung verstanden, welche der aufgehängte Magnetstab in dem betreffenden Augenblick wirklich eben hat, sondern diejenige, welche er haben würde, wenn er (oder genauer zu reden, seine magnetische Axe) in diesem Augenblick genau im magnetischen Meridian wäre. Diese Distinction war unnöthig, so lange man sich nur solcher Nadeln bediente, die eine sehr grosse Genauigkeit nicht geben konnten: man brauchte nur dafür zu sorgen, dass die Nadel um die Zeit der Beobachtung in keiner erkennbaren Schwingung begriffen war, und erhielt damit das Gesuchte unmittelbar. Bei den viel grössern Forderungen, die man an die Genauigkeit der Bestimmungen durch die jetzt eingeführten Apparate machen kann und machen muss, kann aber von einer solchen unmittelbaren Bestimmung nicht mehr die Rede sein. Es steht nicht in unsrer Macht, die Nadel des Magnetometers so vollkommen zu beruhigen, dass gar keine erkennbaren Schwingungsbewegungen zurückbleiben; wenigstens kann es nicht mit Sicherheit ohne Zeitaufwand, und nicht auf die Dauer geschehen. Es werden daher an die Stelle der unmittelbaren Beobachtung solche mittelbare Bestimmungen treten müssen, zu denen eine vollkommene Beruhigung unnöthig ist.

Die sich zuerst darbietende Methode besteht darin, dass man die Nadel absichtlich im schwingenden Zustande beobachtet, zwei auf einander folgende äusserste Stellungen (ein Minimum und ein Maximum) an der Scale aufzeichnet, und zwischen beiden das Mittel nimmt. Dieses an sich unverwerfliche Verfahren erfordert jedoch, wenn die Schwingungen eine beträchtliche Grösse haben, eine Modification, und ist, wenn die Schwingungen klein sind, nur unter einer einschränkenden Bedingung zulässig. Im ersten Fall nemlich wird selbst von einer Schwingung zur andern die successive Abnahme des Schwingungsbogens nicht unmerklich, daher auch schon die Abweichung vom wirklichen Meridian auf der Maximum-Seite geringer sein, als sie beim vorhergehenden Minimum auf der ent\SourceLeaf{549}gegengesetzten Seite gewesen war, folglich das Mittel aus diesem Minimum und dem folgenden Maximum zu klein werden. Aus derselben Ursache wird das Mittel aus diesem Maximum und dem folgenden Minimum ein zu grosses Resultat geben. Da nun aber die Abnahme des Schwingungsbogens einige Schwingungen hindurch beinahe gleichförmig bleibt, so kann man das Mittel aus zwei solchen Mitteln als hinlänglich genau, und zwar als geltend für den Augenblick der zweiten Elongation betrachten. Oder, um es durch eine Formel auszudrücken, wenn $a,b,c$ die Ablesungen in drei auf einander folgenden Elongationen sind (gleich viel, ob die erste und dritte Minima sind, und die zweite ein Maximum, oder umgekehrt), so stellt $\tfrac14(a+2b+c)$ den im Augenblick der Elongation $b$ Statt findenden Stand des magnetischen Meridians dar.

Bei kleinen Schwingungen ist dieses Verfahren nur dann zulässig, wenn die Declination keinen in kurzer Zeit merklichen Veränderungen unterworfen ist, und man kann dann schon das Mittel aus zwei auf einander folgenden Elongationen, als für den in der Mitte liegenden Augenblick gültig ansetzen: im entgegengesetzten Fall aber, d.\ i.\ zu einer Zeit, wo in der Declination schnell beträchtliche Änderungen vorgehen, kann dies Verfahren seine Brauchbarkeit gänzlich verlieren.

Immer aber behält die Methode, den Stand des magnetischen Meridians aus beobachteten Elongationen zu bestimmen, die Unbequemlichkeit, dass die Augenblicke, für welche das erhaltene Resultat gilt, nicht dieselben sind (oder es nur zufällig werden), für welche man den Stand verlangt. Und wenn auch dies in der Mehrzahl der Fälle wenig erheblich sein mag, so verdient doch offenbar ein anderes Verfahren den Vorzug, welches, von jener Inconvenienz frei, Bequemlichkeit, Gleichförmigkeit und alle nur zu wünschende Schärfe in sich vereinigt, und deshalb von sämmtlichen Theilnehmern an den Terminsbeobachtungen befolgt wird.

Dieses Verfahren beruht auf dem Satze, dass das Mittel aus zwei Stellungen der Nadel, die zweien genau um eine Schwingungsdauer von einander abstehenden Augenblicken entsprechen, mit derjenigen Lage des magnetischen Meridians übereinstimmt, welche für das Mittel dieser Zeiten Statt fand, in welche Theile der Schwingungsperiode diese Zeiten auch fallen mögen. Dieser Satz würde in mathematischer Schärfe wahr sein, wenn theils keine äussere Ursachen (wie der Widerstand der Luft u.\ dergl.) zur successiven Verkleinerung des Schwingungs\SourceLeaf{550}bogens. wirkten, theils die etwaige Veränderung in der Lage des magnetischen Meridians während jener kurzen Zwischenzeit nur als \textit{gleichförmig} betrachtet werden dürfte. Der erstere Umstand hat aber gar keinen merklichen Einfluss, wenn man das Verfahren immer nur auf sehr kleine Schwingungsbewegungen anwendet, und was den zweiten betrifft, so sind die Veränderungen der Declination während einer so kurzen Zwischenzeit in der Regel schon an sich kaum merklich, und um so mehr ist man berechtigt, wenigstens die Gleichförmigkeit der Veränderungen während dieser kurzen Zeit gelten zu lassen\HistoricalNote{550}.

Hiemit ist nun die Aufgabe von selbst gelöst. Um den der Declination für die Zeit $T$ entsprechenden Stand der Nadel zu erfahren, braucht man nur, nachdem nöthigenfalls vorher ihre Bewegungen durch angemessene Beruhigungsmittel auf sehr kleine gebracht sind, die wirklichen Stellungen für die Zeiten $T-\tfrac12t$ und $T+\tfrac12t$ zu beobachten, und daraus das Mittel zu nehmen, wo $t$ die Schwingungsdauer bedeutet. Inzwischen, grösserer Genauigkeit und Sicherheit wegen, beschränkt man sich hierauf nicht, sondern macht noch einige ähnliche Bestimmungen für ein Paar Zeitmomente kurz vor, und eben so viele nach $T$, immer in gleichen Intervallen, unter welcher Voraussetzung, insofern während dieser Zeit die Änderung der Declination als gleichförmig betrachtet werden darf, das Mittel aus allen diesen Resultaten das für die Zeit $T$ geltende \textit{Endresultat} sein wird, und zuverlässiger als die einzelne Bestimmung für $T$ selbst.

Die einfachste Art, dies auszuführen, besteht, wenn z.\ B.\ das Endresultat auf fünf partiellen Resultaten beruhen soll, darin, dass man den wirklichen Stand der Nadel für die sechs Zeiten
\[
T-\tfrac52t,\quad T-\tfrac32t,\quad T-\tfrac12t,\quad T+\tfrac12t,\quad T+\tfrac32t,\quad T+\tfrac52t
\]
aufzeichnet. Sind die aufgezeichneten Zahlen $a,b,c,d,e,f$, so wird $\tfrac12(a+b)$ das für die Zeit $T-2t$ geltende Resultat sein; eben so
\[
\tfrac12(b+c),\quad\tfrac12(c+d),\quad\tfrac12(d+e),\quad\tfrac12(e+f)
\]
für die Zeiten $T-t$, $T$, $T+t$, $T+2t$; und das Mittel aus diesen partiellen
\SourceLeaf{551}Resultaten oder der fünfte Theil ihrer Summe wird als berichtigtes Endresultat für die Zeit $T$ anzunehmen sein.

\Needspace{15\baselineskip}
Als Beispiel möge hier das Detail der Beobachtung in Göttingen am 17. August 1836 für $15^{\mathrm h}\,30'$ stehen. Der Beobachter war Hr.\ Dr.\ \textsc{Wappäus}. Für $t$ war angenommen $20''$.
\[
\begin{array}{rrr|r|}
15^{\mathrm h}&29'&10''&865.2\\
&&30&867.5\\
&&50&866.2\\
&30&10&868.0\\
&&30&867.3\\
&&50&868.5
\end{array}
\quad
\left.\begin{array}{r}866.35\\866.85\\867.10\\867.65\\867.90\end{array}\right\}867.16
\]
Die erste Columne enthält hier die Beobachtungszeiten, die zweite die aufgezeichneten Scalentheile, die dritte das Mittel zwischen je zwei auf einander folgenden Aufzeichnungen, mithin die für
\[
15^{\mathrm h}\,29'\,20'',\quad15^{\mathrm h}\,29'\,40'',\quad15^{\mathrm h}\,30'\,0'',\quad15^{\mathrm h}\,30'\,20'',\quad15^{\mathrm h}\,30'\,40''
\]
geltenden partiellen Resultate, und daneben das für $15^{\mathrm h}\,30'\,0''$ geltende Endresultat. In diesem Beispiele ist die im Laufe der Beobachtungen fortwährend Statt habende Veränderung der Declination offenbar, und wird auch durch die vorhergehenden und folgenden Resultate bestätigt. Es war nemlich das Resultat
\[
\begin{array}{lrrll}
\text{für}&15^{\mathrm h}&25'\,0''&.\ .\ .\ .\ .\ .&862.82\\
&&35\quad0&.\ .\ .\ .\ .\ .&872.32
\end{array}
\]
Gewöhnlicher übrigens, als so beträchtliche Änderungen, ist der während der Zeit, welche ein Beobachtungssatz erfordert, fast stationäre Stand der Declination, und in solchen Fällen dient das kleinere oder grössere Hinundherschwanken der partiellen Resultate als ein Maassstab für die grössere oder geringere Zuverlässigkeit der Beobachtungen selbst, möge sie nun von dem Grad der Geschicklichkeit und Aufmerksamkeit des Beobachters, oder der Güte des Apparats selbst, oder von den mehr oder weniger günstigen äussern Umständen abhängen.

Das beschriebene Verfahren ist dasjenige, welches die meisten Theilnehmer an den Terminsbeobachtungen befolgen. Es setzt die Kenntniss der Schwingungsdauer der Nadel voraus, welche bekanntlich zugleich von der Stärke der
